\section{Introduction}
\label{sec:introduction}

% ---------------------------------------------------------
% 1.1 BACKGROUND AND PROBLEM STATEMENT
% ---------------------------------------------------------

\subsection{Background and Problem Statement}
\label{sec:background_problem_statement}

Health insurance cost prediction lies at the intersection of statistical
learning, actuarial modelling and explainable machine learning.  The
empirical task involves estimating a continuous monetary outcome from
policyholder characteristics.  A single fitted value or prediction-error
metric, however, does not fully characterise the behaviour of the fitted
model.  Understanding how predicted costs vary across a simulated
population and identifying which input variables contribute most to high
predicted costs are questions that require methods beyond point
prediction.

Classical actuarial pricing relies on generalised linear models, which
offer transparent coefficient interpretation but are limited to
pre-specified linear predictor structures.  Deep neural networks
provide a flexible alternative, capable of approximating complex
nonlinear relationships between policyholder attributes and cost without
requiring explicit interaction terms
\citep{goodfellow2016deep, hornik1989multilayer}.  Their flexibility,
however, comes at the expense of direct interpretability: the parameters
of a trained network do not admit the same coefficient-level reading as
a linear model.

This tension between predictive flexibility and interpretability
motivates the present study.  A framework is needed that combines
flexible prediction with structured distributional analysis and
principled feature attribution, so that the fitted model can be
examined not only for its accuracy but also for its behaviour under
simulated input variation and for the variables to which it assigns
large contributions in the upper tail of predicted costs.

% ---------------------------------------------------------
% 1.2 POSITIONING WITHIN MATHEMATICAL STATISTICS
% ---------------------------------------------------------

\subsection{Positioning within Mathematical Statistics}
\label{sec:positioning}

This dissertation is positioned within statistical learning,
simulation-based inference, model-implied predictive distributions, and
explainable statistical modelling.  The statistical contribution is not
only the fitting of a predictive model, but the construction of a
model-implied distribution under an explicit input law and the
conditional analysis of feature attributions in the upper tail of that
distribution.

The predictive component draws on supervised regression and
neural-network approximation theory.  The simulation component uses
parametric Monte Carlo methods to propagate input variation through the
fitted prediction function.  The explanation component applies
Shapley-value decomposition, grounded in cooperative game theory, to
attribute fitted predictions to individual input variables.  Each of
these components has a well-established theoretical basis in
mathematical statistics; the contribution lies in their integration
within a single empirical workflow applied to health insurance cost
prediction.

% ---------------------------------------------------------
% 1.3 RESEARCH QUESTIONS
% ---------------------------------------------------------

\subsection{Research Questions}
\label{sec:research_questions}

The dissertation addresses the following research questions:

\begin{enumerate}
    \item How does a deep neural network regression model compare with a
          classical generalised linear model in predicting health insurance
          costs on a benchmark dataset?

    \item What is the shape and distributional behaviour of the
          model-implied predicted-cost distribution obtained by propagating
          parametric Monte Carlo input profiles through the trained network?

    \item Is the product-marginal Monte Carlo input distribution a
          defensible approximation of the training-data joint distribution
          for the purpose of generating simulated profiles?

    \item How sensitive is the predicted-cost distribution to
          residual perturbation from the fitted model?

    \item Which input variables receive the largest Shapley-value
          attributions globally, and how does the feature-importance
          ranking change conditionally within high-cost regions of the
          predicted-cost distribution?

    \item What explains any differences between the global and
          conditional feature-importance rankings in the upper tail?
\end{enumerate}

% ---------------------------------------------------------
% 1.4 RESEARCH OBJECTIVES
% ---------------------------------------------------------

\subsection{Research Objectives}
\label{sec:research_objectives}

The research questions stated above are operationalised through the
following objectives:

\begin{enumerate}
    \item To fit a deep neural network regression model to a health
          insurance benchmark dataset and to evaluate its predictive
          performance against a classical generalised linear model
          benchmark.

    \item To design and implement a parametric Monte Carlo simulation
          that generates synthetic policyholder profiles from fitted
          marginal distributions, propagates them through the trained
          network, and produces a model-implied distribution of
          predicted costs.

    \item To validate the Monte Carlo input distribution against the
          observed training data using joint-distribution diagnostics,
          bootstrap resampling, and a Shapley attribution invariance
          check.

    \item To conduct a residual-augmented Monte Carlo sensitivity
          analysis that examines the effect of residual perturbation on
          the predicted-cost distribution.

    \item To compute Shapley-value attributions for simulated
          policyholder profiles and to summarise feature importance at
          the global level and conditionally within high-cost regions
          defined by upper-tail thresholds.

    \item To examine how the conditional feature-importance ranking in
          high-cost regions differs from the global ranking, and to
          interpret these differences within the context of the fitted
          model and benchmark data.
\end{enumerate}

% ---------------------------------------------------------
% 1.5 RESEARCH GAP
% ---------------------------------------------------------

\subsection{Research Gap}
\label{sec:research_gap_intro}

Existing health-cost prediction studies often emphasise predictive
performance, while interpretability studies often report global feature
importance on observed data only.  The combination of a fitted neural
network, Monte Carlo simulation under an explicit input distribution,
and conditional Shapley summaries in high-cost regions has not been
presented as an integrated workflow in the health insurance literature.
This gap is developed through the reviewed studies in
Section~\ref{sec:literature_review} and stated formally in
Section~\ref{sec:research_gap}.

% ---------------------------------------------------------
% 1.6 METHODOLOGICAL OVERVIEW
% ---------------------------------------------------------

\subsection{Methodological Overview}
\label{sec:methodological_overview}

The empirical analysis follows a three-stage workflow:

\begin{enumerate}
    \item \textbf{Prediction.}  A feed-forward deep neural network is
          trained on a publicly available health insurance benchmark
          dataset using a sequential hyperparameter comparison.  A
          Gamma generalised linear model with a logarithmic link
          function is fitted as a classical actuarial benchmark.

    \item \textbf{Simulation.}  Parametric distributions are fitted to
          the training attributes, and Monte Carlo sampling generates a
          large synthetic policyholder population.  The trained network
          is evaluated at each simulated profile, yielding a
          model-implied predicted-cost distribution.  Convergence
          diagnostics, joint-distribution validation, and
          residual-augmented sensitivity analysis assess the quality and
          robustness of the simulation.

    \item \textbf{Explanation.}  Shapley values are computed for the
          simulated profiles using gradient-based approximations.
          Global feature importance and conditional Shapley summaries
          at the 90th, 95th and 99th percentile thresholds identify the
          variables that drive high predicted costs within the fitted
          model.
\end{enumerate}

The dataset used is the Kaggle \textit{Insurance Data for Machine
Learning} benchmark.  It is treated as a methodological benchmark
rather than a representative underwriting portfolio from a specific
insurer.  All conclusions are internal to the modelling workflow and
the benchmark data.

% ---------------------------------------------------------
% 1.7 SCOPE AND DELIMITATIONS
% ---------------------------------------------------------

\subsection{Scope and Delimitations}
\label{sec:scope_delimitations}

The following delimitations define the boundaries of the study:

\begin{itemize}
    \item \textbf{Benchmark data.}  The analysis uses a publicly
          available synthetic insurance dataset.  The results describe
          relationships learned by the fitted model within the benchmark
          data and should not be interpreted as causal determinants of
          healthcare costs or as direct evidence of real-world pricing
          practice.

    \item \textbf{Limited model comparison.}  A classical generalised
          linear model benchmark is included to contextualise the
          predictive performance of the deep neural network.  However,
          the dissertation does not attempt a comprehensive
          model-selection study across all possible architectures,
          tree-based models, ensembles, or Bayesian alternatives.

    \item \textbf{Single fitted network.}  The Monte Carlo and Shapley
          analyses are applied to a single trained network.  The
          predicted-cost distribution and feature attributions are
          conditional on this particular fitted model and its estimated
          parameters.

    \item \textbf{Independent marginals.}  The Monte Carlo input
          distribution is constructed as a product of independent fitted
          marginals.  This is a deliberate modelling choice that
          broadens the simulated input space but does not reproduce
          the dependence structure of the observed data.

    \item \textbf{Approximate Shapley values.}  The reported
          attributions are gradient-based SHAP approximations, not
          exact game-theoretic Shapley values.  They are conditional on
          the fitted network, the background sample, and the explainer
          settings.

    \item \textbf{No causal inference.}  The Shapley decomposition
          identifies variables to which the fitted model assigns large
          contributions.  It does not establish causal relationships
          between input variables and healthcare expenditure.
\end{itemize}

% ---------------------------------------------------------
% 1.8 DISSERTATION OUTLINE
% ---------------------------------------------------------

\subsection{Dissertation Outline}
\label{sec:dissertation_outline}

The remainder of this dissertation is organised as follows:

\begin{description}
    \item[Chapter~\ref{sec:literature_review}: Literature Review.]
    Reviews the theoretical and empirical foundations of the DNN--Monte
    Carlo--Shapley framework, covering statistical learning, deep neural
    networks for regression, health insurance pricing, Monte Carlo
    simulation, Shapley-value attribution, and the research gap.

    \item[Chapter~\ref{chap:methods_and_data}: Methods and Data.]
    Describes the benchmark dataset, preprocessing pipeline, network
    architecture and training protocol, predictive and structural
    diagnostics, Monte Carlo simulation design, input-distribution
    validation, residual-augmented sensitivity analysis, and
    Shapley-value implementation.

    \item[Chapter~\ref{sec:results}: Results.]
    Presents the hyperparameter comparison results, the generalised
    linear model benchmark comparison, the Monte Carlo predicted-cost
    distribution, the residual-augmented sensitivity analysis,
    Shapley-value attribution at the global and conditional levels,
    and the integration of the Monte Carlo and Shapley findings.

    \item[Chapter~\ref{sec:conclusion}: Conclusion.]
    Summarises the findings, states the contributions and limitations
    of the study, and identifies directions for future research.
\end{description}
